\documentclass[10pt,a4paper]{amsart}
\usepackage[margin=19mm]{geometry}
\usepackage{amsmath,amssymb,mathtools}
\usepackage[scr=rsfs,cal=euler]{mathalfa}
\usepackage{xfrac}
\usepackage[unicode,hidelinks,bookmarks=false]{hyperref}
\newcommand{\ie}{i.e.\ }
\newcommand{\cf}{cf.\ }
\newcommand{\resp}{respectively\ }
\newcommand{\Subsection}{\subsection}
\providecommand{\nobreakdash}{\nobreak\hbox{-}}
\setlength{\emergencystretch}{2em}
\multlinegap0pt
\usepackage{bm}
\usepackage{bbm}
\usepackage{dsfont}
\usepackage{xspace}
\usepackage{cancel}
\usepackage{mathtools}
\usepackage{amsthm}
\makeatletter
\@ifundefined{lemma}{\newtheorem{lemma}{Lemma}}{}
\makeatother
\title[On Empirical Spectral Distributions for Random Tensor Produc]{On Empirical Spectral Distributions for Random Tensor Product Models}
\author{Simona Diaconu}
\date{Source: arXiv:2602.01242v1 (CC BY 4.0)}
\begin{document}
\maketitle

\noindent\emph{Source and licence.} On Empirical Spectral Distributions for Random Tensor Product Models. Version arXiv:2602.01242v1 (CC BY 4.0), distributed under the Creative Commons Attribution 4.0 International licence, as recorded in the arXiv licence metadata for this exact version and shown on the abstract page https://arxiv.org/abs/2602.01242v1. Full rights notice: licenses/arxiv-2602.01242-CC-BY-4.0.txt.\par\smallskip

\noindent\textit{Editorial scope.} Counted unit: the Lemma whose source label is \texttt{lemmamom} in the article (source0.tex lines 1095--1106, source label \texttt{lemmamom}), delivered together with the article's own complete original proof of it (source0.tex lines 1108--1168, one \texttt{proof} environment of 61 lines). No mathematics is written, repaired or re-derived: statement and proof are the article's own text line for line. Required context retained uncounted: none. Every definition, display and label of those lines is retained unchanged. Omitted from this package and not redistributed: 1--1094, 1107--1107, 1169--1227. Nothing inside a retained excerpt is omitted or altered: every retained body is delivered exactly as the article's lines are written, so no omission receipt of any kind accompanies this package. Citations: the retained excerpts carry no citation command at all, so no bibliography entry of the article is transcribed and no reference is redistributed. Reference adaptation: none was needed and none was made; every reference inside the delivered unit resolves inside it, so no unresolved reference or citation is produced. Printed slips kept literal: no printed word, symbol, hypothesis or number of the counted statement or of its proof was altered. Layout and class: the article's own class is \texttt{article} with packages including fontenc, titlesec, inputenc, amsmath, amssymb, amsfonts, amsthm, geometry, babel, xcolor, url, hyperref; the wrapper is the lane's pinned \texttt{amsart} editorial wrapper at 10pt on A4, which restates the theorem-like environments the retained text needs and adds the article's own macro definitions that the retained text needs (none), with extra wrapper packages (bm, bbm, dsfont, xspace, cancel, mathtools). The delivered numbering is the unit's own. Assets: no retained span contains a figure environment, a graphics-inclusion command, a PDF-page inclusion, a table or a TikZ picture; no image file is copied into this package, the package declares no asset dependency and no image file is redistributed with it. Licence: version arXiv:2602.01242v1 of this article is distributed under the Creative Commons Attribution 4.0 International licence, as recorded in the arXiv licence metadata for this exact version and shown on the abstract page https://arxiv.org/abs/2602.01242v1. A full rights notice is delivered at licenses/arxiv-2602.01242-CC-BY-4.0.txt. Characters: every retained excerpt is pure ASCII, so the delivered engine, XeLaTeX with its default Latin Modern text font, reports no missing character.
\begin{lemma}\label{lemmamom}
    Suppose \(\mathbb{E}[x^2]=1,\mathbb{E}[x^{2q}] \leq (Cq)^q\) for all \(q \in \mathbb{N},\) and let 
    \[c(m,n,(d_j)_{1 \leq j \leq m})=\frac{1}{\binom{n}{d_1} \cdot \binom{n}{d_2} \cdot ... \cdot \binom{n}{d_m}}\sum_{1 \leq i_j \leq \binom{n}{d_j}}{\mathbb{E}[Z^2_{d_1,p_1}Z^2_{d_2,p_2}...Z^2_{d_m,p_m}]}\]
    for \(m \in \mathbb{N}, d_1,d_2, \hspace{0.05cm}...\hspace{0.05cm},d_m \in \{1,2,\hspace{0.05cm}...\hspace{0.05cm},n\},\) \(x_1,x_2,\hspace{0.05cm}...\hspace{0.05cm},x_n,\) i.i.d., \(x_1\overset{d}{=} x,\) and \((Z_{l,k})_{1 \leq k \leq \binom{n}{l}}\) a fixed labeling of 
    \((x_{r_1}x_{r_2}...x_{r_l})_{1 \leq r_1<r_2<...<r_l \leq n}.\)
    %for the products of \(l\) pairwise distinct random variables among \(x_1,x_2,\hspace{0.05cm}...\hspace{0.05cm},x_n,\) i.i.d. copies of \(x.\) 
    Then
    \begin{equation}\label{wantedbd}
        c(m,n,(d_j)_{1 \leq j \leq m}) \leq 1+e^{\frac{2C^2e^2(\sum_{1 \leq j \leq m}{d_j})^2}{n}} \cdot \frac{2C^2e^2(\sum_{1 \leq j \leq m}{d_j})^2}{n}
    \end{equation}
    when \(\sum_{1 \leq j \leq m}{d_j} \leq \frac{n^{1/2}}{3Ce}.\)  %K=\max{(Ce,2C^2e^2)} 
\end{lemma}

\begin{proof}
    %The case \(m=1\) is immediate from independence and \(\mathbb{E}[x^2]=1.\) Suppose next \(m \geq 2.\) 
    If there is no overlap among the positions underlying \(p_1,p_2,\hspace{0.05cm}...\hspace{0.05cm},p_m,\) then the expectations contribute at most \(1\) from independence and \(\mathbb{E}[Z^2_{d_j,p_j}]=1\) for \(1 \leq j \leq m.\) Else, there is an index \(r \in \{1,2,\hspace{0.05cm}...\hspace{0.05cm},n\}\) that appears exactly in the sets underpinning the positions \(p_{i_1},p_{i_2},\hspace{0.05cm}...\hspace{0.05cm},p_{i_k}\) for \(1 \leq i_1<i_2<...<i_k \leq m\) and some \(2 \leq k \leq m.\) The contribution of \(x_r\) to the expectation is, by independence, \(\mathbb{E}[x^{2k}],\) and
    \[\binom{n}{d}=\binom{n-1}{d-1} \cdot \frac{n}{d} \hspace{0.5cm} (1 \leq d \leq n), \hspace{0.8cm} \binom{n}{d}=\binom{n-1}{d} \cdot \frac{n}{n-d} \hspace{0.5cm} (1 \leq d \leq n-1)\]
    imply
    \begin{equation}\label{ineqq1}
        c(m,n,(d_j)_{1 \leq j \leq m}) \leq 1+\sum_{2 \leq k \leq m}{n \cdot \mathbb{E}[x^{2k}]\sum_{1 \leq i_1<...<i_k \leq m}{\frac{1}{\frac{n}{d_{i_1}} \cdot \frac{n}{d_{i_2}} \cdot ... \cdot \frac{n}{d_{i_k}}}} \cdot c(m,n-1,(d_j-\chi_{j \in \{i_1,i_2, \hspace{0.05cm}...\hspace{0.05cm},i_k\}})_{1 \leq j \leq m})}.
    \end{equation}
    The multinomial theorem gives 
    \[\sum_{1 \leq i_1<...<i_k \leq m}{\frac{1}{\frac{n}{d_{i_1}} \cdot \frac{n}{d_{i_2}} \cdot ... \cdot \frac{n}{d_{i_k}}}} \leq (\sum_{1 \leq j \leq m}{\frac{1}{\frac{n}{d_j}}})^k \cdot \frac{1}{k!}=(\frac{\sum_{1 \leq j \leq m}{d_j}}{n})^k \cdot \frac{1}{k!},\] 
    from which (\ref{ineqq1}) can be changed to
    \[c(m,n,(d_j)_{1 \leq j \leq m}) \leq 1+\sum_{2 \leq k \leq m}{n \cdot \mathbb{E}[x^{2k}] \cdot (\frac{\sum_{1 \leq j \leq m}{d_j}}{n})^k \cdot \frac{1}{k!} \cdot \max_{1 \leq i_1<...<i_k \leq m}}{c(m,n-1,(d_j-\chi_{j \in \{i_1,i_2, \hspace{0.05cm}...\hspace{0.05cm},i_k\}})_{1 \leq j \leq m})}.\]
    \par
    Proceed by induction on \(s=\sum_{1 \leq j \leq m}{d_j} \leq \frac{n^{1/2}}{3Ce}.\) Note that \(n \geq 2\) from \(s \geq 1,3Ce \geq 3e>1\) as
    \begin{equation}\label{ceq}
        C \geq \mathbb{E}[x^2]=1.
    \end{equation}
    The base case \(s=1\) has \(m=1,\) and the result is clear (the sum is at most \(1\)). Suppose next the bound holds for \(s-1 \geq 1.\) By virtue of \(k! \geq (\frac{k}{e})^k,\) it suffices to show that\footnote{The bound in (\ref{wantedbd}) does  not depend directly on \(m,\) and thus the cases in which \(d_{i_j}=1\) for some \(1 \leq j \leq k\) cause no issue (such positions can be dropped from the tuple corresponding to \(n-1\)).}
    \begin{equation}\label{eqsum}
        1+\sum_{2 \leq k \leq m}{n \cdot (Ck)^{k} \cdot (\frac{es}{nk})^k \cdot (1+e^{\frac{2C^2e^2(s-k)^2}{n-1}} \cdot \frac{2C^2e^2(s-k)^2}{n-1})}:=1+I+II
    \end{equation}
    is upper bounded by the right-hand side term of (\ref{wantedbd}) (using as well that \(x \to xe^{x}\) increases on \([0,\infty)\)), where \(I,II\) are obtained by opening the parentheses in the last factors of the summands since 
     \[s-k \leq \frac{n^{1/2}}{3Ce}-2 \leq \frac{(n-1)^{1/2}}{3Ce}\] 
     from 
     \[\frac{n^{1/2}-(n-1)^{1/2}}{3Ce}=\frac{1}{3Ce \cdot (n^{1/2}+(n-1)^{1/2})} \leq \frac{1}{3Ce} \leq \frac{1}{3e}<1\]
     via (\ref{ceq}). 
     \par
     Rewrite the factors in (\ref{eqsum}) as
    \[n \cdot (Ck)^{k} \cdot (\frac{es}{nk})^k=n \cdot (\frac{Ces}{n})^{k}=\frac{C^2e^2s^2}{n} \cdot (\frac{Ces}{n})^{k-2},\]
    from which it is enough to show that
    \[I=\frac{C^2e^2s^2}{n} \sum_{2 \leq k \leq m}{(\frac{Ces}{n})^{k-2}} \leq e^{\frac{2C^2e^2s^2}{n}} \cdot \frac{C^2e^2s^2}{n},\]
    \[II=\frac{C^2e^2s^2}{n} \sum_{2 \leq k \leq m}{(\frac{Ces}{n})^{k-2} \cdot e^{\frac{2C^2e^2(s-k)^2}{n-1}} \cdot \frac{2C^2e^2(s-k)^2}{n-1}} \leq e^{\frac{2C^2e^2s^2}{n}} \cdot \frac{C^2e^2s^2}{n}\]
    when \(m \leq s \leq \frac{n^{1/2}}{3Ce}\) to complete the induction step. %(these inequalities are applied to \(s=\sum_{1 \leq j \leq m}{d_j} \geq \sum_{1 \leq j \leq m}{1}=m).\)
    \par
    Begin with \(I:\) the desired claim is equivalent to
    \[\sum_{0 \leq k \leq m-2}{(\frac{Ces}{n})^{k}} \leq e^{\frac{2C^2e^2s^2}{n}},\]
    and ensues from
    \begin{equation}\label{helper}
        \sum_{0 \leq k \leq M}{x^k} \leq e^{Mx} \hspace{1cm} (x \geq 0,M \geq 0),
    \end{equation}
    a consequence of
    \[e^{Mx}=\sum_{k \geq 0}{\frac{(Mx)^k}{k!}} \geq \sum_{k \geq 0}{\frac{(Mx)^k}{k^k}} \geq \sum_{0 \leq k \leq M}{x^k},\]
    under the convention \(0^0:=1.\) Namely, (\ref{helper}) gives
    \[\sum_{0 \leq k \leq m-2}{(\frac{Ces}{n})^{k}} \leq e^{\max{(m-2,0)} \cdot \frac{Ces}{n}}<e^{s \cdot \frac{Ces}{n}}< e^{\frac{2C^2e^2s^2}{n}}\]
    by using \(s \geq m\) and (\ref{ceq}). %2C \geq \mathbb{E}[x^2]=1.\)
    \par
    Consider now \(II,\) for which it suffices to justify
    \[\frac{C^2e^2s^2}{n}\sum_{2 \leq k \leq m}{(\frac{Ces}{n})^{k-2} \cdot e^{\frac{2C^2e^2(s-k)^2}{n-1}}} \leq e^{\frac{2C^2e^2s^2}{n}} \cdot \frac{n-1}{2n}\]
    from \((s-k)^2 \leq s^2\) for \(k \leq m \leq s.\) Since \((s-k)^2 \leq s(s-k),\) and the corresponding terms, after replacing the squares in the exponentials by these linear bounds in \(k,\) satisfy
    \[\frac{t(k+1)}{t(k)}=\frac{Ces}{n} \cdot e^{-\frac{2C^2e^2s}{n-1}}< \frac{n^{1/2}}{3n} \leq \frac{1}{3}.\]
    This renders that the sum is at most \(\frac{1}{1-\frac{1}{3}} \cdot t(2) \leq 2t(2),\) amounting to an overall bound of at most
    \[\frac{C^2e^2s^2}{n} \cdot 2 \cdot e^{\frac{2C^2e^2s(s-2)}{n-1}} \leq e^{\frac{2C^2e^2s^2}{n}} \cdot \frac{n-1}{2n}:\]
    to derive this last inequality, use
    \[e^{\frac{2C^2e^2s(s-2)}{n-1}}< e^{\frac{2C^2e^2s^2}{n}}\]
    from 
    \[\frac{s-2}{n-1}-\frac{s}{n}=\frac{n(s-2)-s(n-1)}{n(n-1)}=\frac{s-2n}{n(n-1)} \leq \frac{-n}{n(n-1)}<0\] 
    as \(s \leq \frac{n^{1/2}}{3Ce} \leq \frac{n^{1/2}}{3e}<n,\) while 
    \[\frac{C^2e^2s^2}{n} \cdot 2 \cdot \frac{2n}{n-1} \leq C^2e^2 \cdot \frac{1}{9C^2e^2}\cdot 2 \cdot 4=\frac{8}{9}<1,\]
    completing the induction step (recall that \(n \geq 2\)).
    
\end{proof}

\end{document}
